% Source: https://artofproblemsolving.com/community/c3h1458639_centroids
% Stable identity: aops-topic-1458639-q01
Let $M$ be the midpoint of $BC$; the median $AM$ passes through $O$.

Triangles $ABM$ and $ACM$ have equal bases $BM=MC$ and the same height from $A$, so $[ABM]=[ACM]$. For the same reason $[OBM]=[OCM]$. Subtracting,
\[
[AOB]=[ABM]-[OBM]=[ACM]-[OCM]=[AOC].
\]
Using the median from $B$ (through the midpoint of $CA$) in the same way gives $[BOA]=[BOC]$. So the three triangles $AOB$, $BOC$, $COA$ have equal areas. They do not overlap and together they fill triangle $ABC$ (since $O$ is inside it), so each has area $\dfrac{1}{3}[ABC]$. $\blacksquare$

\begin{asy}
usepackage("amsmath");
size(6cm);
defaultpen(fontsize(10pt));
pair A=(1.4,4),B=(0,0),C=(6,0);
pair G=(A+B+C)/3, M=(B+C)/2, K=(A+C)/2, L=(A+B)/2;
fill(A--B--G--cycle,rgb(0.85,0.9,1)); fill(B--C--G--cycle,rgb(1,0.88,0.85)); fill(C--A--G--cycle,rgb(0.88,1,0.88));
draw(A--M,gray+dashed); draw(B--K,gray+dashed); draw(C--L,gray+dashed);
draw(A--G,blue); draw(B--G,blue); draw(C--G,blue);
draw(A--B--C--cycle,linewidth(0.9));
dot(A);dot(B);dot(C);dot(G,red);dot(M);
label("$A$",A,N); label("$B$",B,SW); label("$C$",C,SE); label("$O$",G,1.3*dir(65),red); label("$M$",M,S);
label("\mbox{}",(min(currentpicture,user=true)+max(currentpicture,user=true))/2,black); // keeps the text after the figure black when LaTeX typesets the labels inline
\end{asy}
